% ====== Preambulo compartido — tesis TCLab (LQR + RL) ======
% Estandar definido en .claude/rules/thesis-format.md

% --- Idioma (primero: condiciona el resto) ---
\usepackage[spanish,es-tabla,es-noquoting]{babel}
\usepackage{csquotes}

% --- Diseno de pagina ---
\usepackage[left=3cm,right=2.5cm,top=2.5cm,bottom=2.5cm]{geometry}
\usepackage{setspace}
\onehalfspacing
\usepackage{fancyhdr}
\pagestyle{fancy}
\fancyhf{}
\fancyfoot[C]{\thepage}
\renewcommand{\headrulewidth}{0pt}

% --- Tipografia (XeLaTeX) ---
\usepackage{fontspec}
\setmainfont{Latin Modern Roman}
\usepackage{microtype}

% --- Titulos ---
\usepackage{titlesec}
\usepackage[title]{appendix}

% --- Matematicas ---
\usepackage{amssymb, amsmath, amsfonts, mathtools}
\usepackage{amsthm}
\theoremstyle{plain}
\newtheorem{teorema}{Teorema}[chapter]
\newtheorem{proposicion}[teorema]{Proposición}
\newtheorem{lema}[teorema]{Lema}
\newtheorem{corolario}[teorema]{Corolario}
\theoremstyle{definition}
\newtheorem{definicion}[teorema]{Definición}
\newtheorem{supuesto}[teorema]{Supuesto}
\DeclareMathOperator*{\argmax}{arg\,max}
\DeclareMathOperator*{\argmin}{arg\,min}
\DeclareMathOperator{\tr}{tr}
\newcommand{\norm}[1]{\left\lVert #1 \right\rVert}
\newcommand{\transpose}{^{\mathsf{T}}}

% --- Tablas ---
\usepackage{array, booktabs, makecell}
\usepackage{siunitx}
\sisetup{output-decimal-marker={,}}
\usepackage[flushleft]{threeparttable}
\usepackage{tabularx, rotating}
\usepackage{tabularray}
\UseTblrLibrary{booktabs, siunitx}

% --- Figuras ---
\usepackage{graphicx, subcaption}
\usepackage{float}
\usepackage{caption}
\captionsetup{font=small, labelfont=bf, justification=justified}

% --- Diagramas de bloques ---
\usepackage{tikz}
\usetikzlibrary{shapes, arrows.meta, positioning, calc}

% --- Codigo fuente ---
\usepackage{listings}
\lstset{basicstyle=\ttfamily\small, breaklines=true, frame=single,
        showstringspaces=false, captionpos=b}

% --- Listas ---
\usepackage{enumitem}

% --- Bibliografia: biblatex + biber, estilo IEEE ---
\usepackage{xurl}
\usepackage{xcolor}
\definecolor{citationcolor}{RGB}{0, 90, 160}
\usepackage[backend=biber, style=ieee, sorting=none, maxbibnames=99]{biblatex}
\addbibresource{Bibliography_base.bib}

% --- Notas al pie ---
\interfootnotelinepenalty=10000

% --- hyperref (penultimo) ---
\usepackage[breaklinks, colorlinks=true,
            linkcolor=citationcolor, citecolor=citationcolor,
            urlcolor=citationcolor]{hyperref}

% --- cleveref (despues de hyperref) ---
\usepackage[spanish,nameinlink]{cleveref}
